%% kapa: exam=1 round=32 year=2021 period=1 form=A subject=civil_law subject_ko=민법 seq=3 number=3
%% source: Q-Net 감정평가사 기출문제 (cst00309) artlSeq=5208397 file=extracted/1008/5208397/A형/2021년도 제32회 감정평가사 1차 1교시 A형.pdf
%% conversion: kordoc 4.24.51 -> tools/gen1.py (auto); review=proofread-sonnet; flags=none
\begin{question}{3}
제한능력에 관한 설명으로 옳지 않은 것은?
\begin{choices}
  \item 가정법원은 한정후견개시의 심판을 할 때 본인의 의사를 고려하지 않아도 된다.
  \item 가정법원은 취소할 수 없는 피성년후견인의 법률행위의 범위를 정할 수 있으나, 성년후견인의 청구에 의하여 이를 변경할 수 있다.
  \item 성년후견인은 일상생활에 필요하고 그 대가가 과도하지 않은 피성년후견인의 법률행위를 취소할 수 없다.
  \item 가정법원은 성년후견개시의 심판을 할 때 본인의 의사를 고려하여야 한다.
  \item 피성년후견인이 성년후견인의 동의를 얻어 재산상의 법률행위를 한 경우에도 성년후견인은 이를 취소할 수 있다.
\end{choices}
\end{question}
\answer{1}
